\section{The bounded ternary block counts}\label{sec:bounded-ternary}
Put $\gamma=(\log_2 3)/3<17/32$. All fibres in this section use the
complete source $J\leq S_b$, without any restriction on its orbits.

\begin{lemma}[Finite multiplicative comparison]\label{lem:finite-multiplicative}
Fix finitely many actual actions $U$ of widths $w$, and fixed finite
comparators $Q_U$ with faithful actions of degrees $0<v_U<h(w)$.
Suppose
\[
 Z_J(U)\leq C_U2^{\eta_Ub}Z_J(Q_U),\qquad
 \delta_U=(h(w)-v_U)/8-\eta_U>0.
\]
Their complete pointed union has an exponentially small forward row
and a quadratically small scalar, with original normalizers $a(U)$.
The comparators need not be quotients of $U$.
\end{lemma}
\begin{proof}
Set $\alpha=v/8+\delta/2$, and call $J$ hot when
$Z_J(Q)>2^{\alpha b}$. The cold coefficient is exactly
$\Delta(n,b)C_U2^{(\eta_U+\alpha)b}/a(U)$.
Its exponent is $-\delta b/2-w^2/16+O_w(\log(n+2))$.
For the full hot fibre use Lemma~\ref{lem:quotient-moment} with
$q=\max(1,\lceil4\delta b/v^2\rceil)$ and $q_*=4\delta b/v^2$.
The normalized leading exponent is
\[
 \frac{(b+qv)^2}{16}-(q-1)\alpha b+\eta b
             -\frac{(b+w)^2}{16}
 =-\frac{\delta^2b^2}{v^2}-\frac{\delta b}{2}
  -\frac{(w-h(w))b}{8}-\frac{w^2}{16}
                  +\frac{v^2(q-q_*)^2}{16}.
\]
The rounding term is bounded and the coarse counting error is
$O(n^{3/2}+n\log(n+2))$. The original hot weight is
$(b!a(U)p_nG_{\lfloor n/2\rfloor})^{-1}$ times the complete hot
sum of $Z_J(U)$. A finite union therefore preserves both strict
losses. Choose certificates before the any-hot/cold partition.
\end{proof}

\begin{lemma}[Central binary comparison]\label{lem:central-comparator}
If $X$ has central subgroup $C_2$ and quotient $T$, then
\[
 Z_J(X)\leq\nu(X)|\operatorname{Hom}(J,C_2)|Z_J(T),\qquad
 |\operatorname{Hom}(J,C_2)|Z_J(T)
       \leq2^{d_2(T)}Z_J(T\times C_2),
\]
where $\nu(X)$ is its number of literal normal subgroups.
\end{lemma}
\begin{proof}
For each $N\lhd X$, $X/N$ is a central extension of a quotient of
$T$ by a group of order at most two. A fixed map to that quotient has
no lifts or at most $|\operatorname{Hom}(J,C_2)|$ lifts. Summing
over all $N$ proves the first inequality, including repetitions of
the same normal subgroup of $T$.
For the second put $h=|\operatorname{Hom}(J,C_2)|$. For each
$S\lhd T$ put $Q=T/S$, $a=|\operatorname{Hom}(Q,C_2)|$.
An onto map $J\to Q$ and a binary character are jointly onto
$Q\times C_2$ precisely when that character does not factor through
the map. Thus the distinct product normals $S\times1,S\times C_2$
contribute $(h-a+1)|\Epi(J,Q)|$. Whenever the latter epimorphism
set is nonempty, $h\geq a$ and
$h\leq a(h-a+1)\leq2^{d_2(T)}(h-a+1)$.
All remaining product normals add nonnegative terms.
\end{proof}

\begin{proposition}[Regular ternary exceptions]\label{prop:bounded-c3}
Every actual transitive action with exact local regular $C_3$ and
two or six blocks has a complete exponentially small forward row
and quadratically small scalar.
\end{proposition}
\begin{proof}
Write $U\leq C_3\wr T$, $E=U\cap C_3^s$. At $s=6$, every
transitive $T\leq S_6$ contains a semiregular $C_3$: a Sylow
three-subgroup of order three has trivial point stabilizers, while
one of order nine is a Sylow subgroup of $S_6$ and contains two
disjoint three-cycles. The ambient module and its dual restrict to
two regular $\F_3[C_3]$ modules. Since
$\F_3[C_3]\cong\F_3[t]/(t^3)$, every submodule has at most two
generators: its number of Jordan blocks is its socle dimension,
at most two. Quotients retain this bound.
The full-chief proof of Theorem~\ref{thm:affine-chief}, with this
capacity, gives $Z_J(U)\leq C_U2^{2\gamma b}Z_J(T)$.
Its physical width is eighteen, its comparator degree six, and
$3/2-2\gamma>0$, so Lemma~\ref{lem:finite-multiplicative} applies.

At $s=2$, exact local $C_3$ forces $E$ to project onto each
coordinate. Its only $C_2$-invariant possibilities in $\F_3^2$ are
the diagonal, the antidiagonal, and the full plane. Coprimality
splits the extension, giving $C_6,S_3,C_3\times S_3$, respectively,
all in their original degree-six actions. The first has full fibre
$|\operatorname{Hom}(J,C_6)|\leq3^{b/3}$. For the other two use
the common source $F=J'J^2$, $B=J/F$. The semisimple module
$H_1(F,\F_3)$ has trivial multiplicity $d=d_3(J)$ and nontrivial
sign multiplicities $a_\theta$, with
$d+\sum a_\theta\leq b/3$. Coprime inflation--restriction gives
\[
 |\Epi(J,S_3)|=3\sum_{\theta\ne0}(3^{a_\theta}-1),\qquad
 |\Epi(J,C_3\times S_3)|
   =(3^d-1)3\sum_{\theta\ne0}(3^{a_\theta}-1).
\]
The second joint map is onto because $C_3,S_3$ have no common
nontrivial quotient. Each expression is at most $3^{b/3+1}$.
The complete normal quotients of $C_3\times S_3$ are
$1,C_2,C_3,C_6,S_3,C_3\times S_3$; the cyclic terms have the same
abelianization bound. Thus $Z_J(U)\leq C_U2^{\gamma b}$.
Its slope is less than $h(6)/8=3/4$, so direct normalization
already gives an exponentially small row and zero scalar.
\end{proof}

\subsection{Signed modules on a bounded number of triples}
For exact natural $S_3$ on $s$ blocks write
\[
 E=U\cap C_3^s,\quad R=U/E\leq C_2\wr T,\quad
 D=\ker(R\longrightarrow T)\leq C_2^s.
\]
These are the actual kernel and actual top of the chosen block
system. If $E=1$, the whole block kernel is trivial: a nontrivial
normal local image is $C_3$ or $S_3$, and in the latter case its
derived subgroup projects onto $C_3$. Hence $U\cong T$ has faithful
degree $s$ and satisfies Theorem~\ref{thm:finite-comparator}.
If an ambient signed ternary module and its dual have capacity $H_3$,
the full-chief argument gives
\begin{equation}\label{eq:bounded-s3-layer}
 Z_J(U)\leq C_U2^{\gamma H_3b}Z_J(R).
\end{equation}
Adding a binary capacity $H_2$ gives the same bound with
$2^{(\gamma H_3+H_2/2)b}Z_J(T)$. These bounds include all proper
kernels, normal graphs and nonsplit extensions.

\begin{lemma}[Signed isotypes]\label{lem:signed-isotypes}
Assume $D\ne1$. Its coordinate characters are all nonzero and have
$q$ distinct values, each repeated $a=s/q$ times. They span $D^*$;
thus $q\mid s$, $d=\dim D\leq q\leq2^d-1$, and $q=1,2$ forces
$d=1,2$, respectively. If $a$ is a power of two or $a=3$, every
submodule of the signed ambient module and of its dual is cyclic.
\end{lemma}
\begin{proof}
Transitivity makes the character multiplicities equal, and their
common kernel is zero because $D$ is its actual coordinate image.
The idempotents of $\F_3[D]$ decompose a submodule into isotypes.
The stabilizer of one isotype acts transitively on its $a$ coordinates
by signed permutations. If $a$ is a power of two, a transitive Sylow
two-subgroup of its coordinate image lifts to a two-group of signed
permutations. The coordinate module is cyclic and semisimple over
this group; an equivariant projection of a cyclic generator generates
any submodule. If $a=3$, a Sylow three-subgroup maps isomorphically
onto a coordinate three-cycle. Its signs have product one, so it
untwists to the regular $\F_3[C_3]$ module, whose ideals are cyclic.
Transport a generator of the chosen isotype by $R$ to generate all
isotypes. Signed permutation matrices identify the dual with the
original module. Quotients of cyclic modules are cyclic as well.
\end{proof}

\begin{lemma}[Two common-source onto bounds]\label{lem:cyclic-dual-onto}
(i) Let $V$ be elementary abelian of order $3^e$, and $R$ a
two-group acting faithfully with cyclic dual $V^*$. Then
\[
 |\Epi(J,V\rtimes R)|\leq3^e|\mathrm{GL}_e(3)|3^{b/3}.
\]
The same conclusion holds for $R=C_6$ acting faithfully on a single
unipotent ternary block, with its involution acting as $-1$.

(ii) If $X$ has a normal Sylow $p$-subgroup $P$ with cyclic center
and $C_X(P)=Z(P)$, then
\[
 |\Epi(J,X)|\leq |Z(P)|\,|\operatorname{Aut}(P)|p^{b/(p-1)}.
\]
\end{lemma}
\begin{proof}
For (i), fix $F=O^2(J)$ in the two-group case and $F=J'$ in the
$C_6$ case. Formation residuals commute with onto maps: the image
of the source residual has the required quotient type, and the
preimage of the target residual has that type, proving both
inclusions. Consequently every map under consideration sends this
same $F$ onto $V$. Put $B=J/F$, $A=H_1(F,\F_3)$.
The restriction is an onto $B$-map $A\to V_\theta$ and dualizes
to an injection $\iota:V_\theta^*\to A^*$. Fix a cyclic vector
$v_0\in V^*$. Its image $f$ determines the whole submodule
$\langle Bf\rangle$. For each $f$, at most
$|\mathrm{GL}_e(3)|$ identifications with $V^*$ are possible, and
faithfulness determines the quotient map $\theta$ uniquely from
each identification. Thus all quotient maps and onto restrictions
together have at most $|\mathrm{GL}_e(3)|3^{\dim A}$ choices.
For fixed restriction the lift fibre is empty or a cocycle torsor
on $B$. Averaging gives $H^1(B,V_\theta)=0$ when $B$ is a
two-group. In the second case $B$ is abelian and its two-primary
Hall subgroup has zero invariants because it maps onto the scalar
involution; averaging and inflation--restriction give the same
conclusion. The torsor has size at most $|V|$. Finally
$\dim A=d_3(F)\leq b/3$.

For (ii), fix $F=O^{p'}(J)$ and its largest $p$-group quotient
$Q=F/O^p(F)$. Every onto map to $X$ sends $F$ onto $P$ and factors
there through $Q$. Since $F\leq S_b$,
$|Q|\leq p^{v_p(b!)}\leq p^{b/(p-1)}$.
A faithful linear character of $Z(P)$ induces an irreducible
constituent faithful on $P$: every nontrivial normal subgroup of a
$p$-group meets its center. Fix its faithful character $\chi$.
Pullback along onto maps $Q\to P$ gives irreducible characters of
$Q$; one pulled-back character determines the kernel and therefore
at most $|\operatorname{Aut}(P)|$ maps. The degree-square identity
bounds their number by $|\operatorname{Aut}(P)||Q|$.
For a fixed onto restriction, two extensions to $J$ differ in
$C_X(P)=Z(P)$, so their quotient maps coincide. Their differences
form a cocycle torsor on the $p'$-group $J/F$ with coefficients
$Z(P)$. Coprime averaging bounds its size by $|Z(P)|$.
\end{proof}

\begin{proposition}[Two symmetric triples]\label{prop:two-s3}
All exact-local-$S_3$ actions on two blocks satisfy a complete finite
forward estimate.
\end{proposition}
\begin{proof}
The block stabilizer is its kernel, so its binary image is either
$C_2^2$ or diagonal $C_2$. In the first case distinct signs force
$E=\F_3^2$ and $U=S_3\wr C_2$. The ternary kernel is faithful
irreducible for $R=D_8$, and Lemma~\ref{lem:derived-affine} gives
$Z_J(U)\leq Z_J(D_8)+216\,3^{b/6}$. The comparator degree is four.
In the diagonal case $R$ is $C_4$ or $V_4$. Coprimality splits
the extension. For $C_4$, its square is $-I$, so $E=\F_3^2$ is
irreducible and the derived-source bound gives
$Z_J(U)\leq Z_J(C_4)+18\,3^{b/3}\leq19\,3^{b/3}$.
For $V_4$, the ambient module has two distinct sign lines; choosing
one or both gives $S_3\times C_2$ or $S_3^2$.
The former has a faithful degree-five action.

For the latter use the same $F=J'J^2$ and sign multiplicities
$a_\theta$ as above. An ordered pair of onto maps to $S_3$ is onto
$S_3^2$ precisely when its binary characters differ: a proper
subdirect has common quotient $C_2$ or $S_3$, forcing equality.
Hence
\[
 |\Epi(J,S_3^2)|=9\sum_{\theta\ne\psi}
 (3^{a_\theta}-1)(3^{a_\psi}-1)\leq18\,3^{b/3}.
\]
The last inequality expands $\prod_\theta(1+(3^{a_\theta}-1))$
and uses $\sum a_\theta\leq b/3$. Its complete normal menu is
\[
 Z_J(S_3^2)=|\Epi(J,S_3^2)|+2|\Epi(J,S_3\times C_2)|
 +2|\Epi(J,S_3)|+|\Epi(J,V_4)|+3|\Epi(J,C_2)|+1.
\]
Indeed intersection with $C_3^2$ is zero, one of its two sign
lines, or the full plane; self-centralization resolves the zero
case, and the other cases give exactly these quotients. Thus
$Z_J(S_3^2)\leq2Z_J(S_3\times C_2)+18\,3^{b/3}$.
Every slope is below $h(6)/8$ and every comparator degree is below
six, so Theorem~\ref{thm:finite-comparator} finishes all cases.
\end{proof}

\begin{proposition}[Four blocks with two signs]\label{prop:four-s3-two-signs}
All exact-local-$S_3$ four-block actions whose binary kernel has two
distinct coordinate characters have a complete finite forward estimate.
\end{proposition}
\begin{proof}
Their equal character classes have size two, $D=C_2^2$, and the
top preserves the two-class partition. Thus $T\leq D_8$ has order
four or eight, and $R$ is a two-group of order sixteen or thirty-two.
Coprimality splits $U=E\rtimes R$. The two isotypes have equal
dimension in every submodule, so nonzero $E$ has dimension two or
four. Every two-dimensional section has monomial image exactly
$D_8\leq\mathrm{GL}_2(3)$: its diagonal signs are independent
and its two lines are interchanged. Its action kernel $K$ in $R$
has order two or four. Such a normal kernel has a central series
of at most two order-two factors, first by meeting $Z(R)$ and
then because a normal subgroup of order two is central.

Put $P=\F_3^2\rtimes D_8=S_3\wr C_2$ in its faithful degree-six
action and $Q_*=P\times C_2^2$ in degree ten. The central series,
Lemma~\ref{lem:central-comparator}, and the literal quotient-menu
inclusion imply $Z_J(V\rtimes R)\leq A_VZ_J(Q_*)$ for every
two-dimensional section $V$, with a fixed finite $A_V$.
They also give $Z_J(R)\leq A_RZ_J(Q_*)$ whenever such a section
exists. No splitting of the central extensions is required.

If $E$ is the full four-dimensional module, it is faithful and its
dual is cyclic: the coordinate module is cyclic and semisimple for
the two-group $R$. Lemma~\ref{lem:cyclic-dual-onto}(i) bounds its
onto maps by $C_0 3^{b/3}$, where
$C_0=81|\mathrm{GL}_4(3)|$. A normal subgroup disjoint from $E$
centralizes $E$, so is trivial. All other proper intersections are
two-dimensional sections. Consequently
\[
 Z_J(U)\leq |\Epi(J,U)|+Z_J(R)
       +\sum_{\substack{I<E\\\dim I=2, I\ R\text{-stable}}}Z_J(U/I).
\]
If $E$ is reducible, the preceding comparisons bound this by
$B_UZ_J(Q_*)+C_0 3^{b/3}$. If it is irreducible, the sum is empty
and use $R$ in faithful degree eight instead. If $\dim E=2$, the
central comparison already bounds the entire fibre. These are all
normal strata. Comparator degrees at most ten and slope $\gamma$
are strictly below the required width-twelve thresholds.
\end{proof}

\begin{proposition}[Three blocks with one sign]\label{prop:three-s3-one-sign}
All exact-local-$S_3$ three-block actions with diagonal binary kernel
have a complete finite forward estimate.
\end{proposition}
\begin{proof}
The top is $C_3$ or $S_3$. A lifted coordinate three-cycle of order
three untwists its signs. In the first case $R=C_2\times C_3$.
In the second the determinant is nontrivial on the scalar binary
kernel and splits $R\to S_3$; a transposition conjugating the
three-cycle to its inverse has constant signs. Thus after diagonal
rescaling $R=C_2\times S_3$ acts by scalar signs and ordinary
coordinate permutations. Its nonzero invariant ternary submodules
are exactly the three ideals in $\F_3[t]/(t^3)$, of lengths
$e=1,2,3$. All are stable under $S_3$ as well.
The scalar involution has zero invariants on the module and on
every quotient. Averaging and inflation--restriction give
$H^1(R,\F_3^3/E)=0$; a translation therefore removes the cocycle
of the actual embedded subgroup. Hence $U_e=E_e\rtimes R$.

For top $C_3$, $U_1=S_3\times C_3=:Q$, faithfully of degree six.
At $e=2,3$, the action is faithful, $U_e'=E_e$, and the dual is
cyclic. Lemma~\ref{lem:cyclic-dual-onto}(i) gives
$|\Epi(J,U_e)|\leq C_e3^{b/3}$.
Normal intersections with $E_e$ form its uniserial chain. At
surviving length at least two, faithfulness forces a normal subgroup
disjoint from the ternary kernel to be trivial. At length one the
quotient is $Q$, whose normals disjoint from its first ternary line
are $1$ and its central $C_3$. There are no mixed normal graphs,
because the involution inverts one line and fixes the other.
Thus, exactly,
\begin{align*}
 Z_J(U_1)&=Z_J(Q),\\
 Z_J(U_2)&=Z_J(Q)+|\Epi(J,U_2)|,\\
 Z_J(U_3)&=Z_J(Q)+|\Epi(J,U_2)|+|\Epi(J,U_3)|.
\end{align*}

For top $S_3$, $U_1=S_3^2=:Q_*$, again faithfully of degree six.
At lengths two and three the action of $R$ is faithful. The normal
Sylow three-subgroup $P_e=E_e\rtimes C_3$ has center the first
ideal $C$, of order three, and $C_{U_e}(P_e)=C$. This follows
since centralizing $E_e$ forces an element into $E_e$, where
centralizing the top three-cycle forces it into $C$.
Lemma~\ref{lem:cyclic-dual-onto}(ii) gives
$|\Epi(J,U_e)|\leq3|\operatorname{Aut}(P_e)|3^{b/2}$.
Every nontrivial normal subgroup contains $C$, by the uniserial
chain and faithfulness. Moreover $U_2/C\cong Q_*$ and
$U_3/C\cong U_2$. To see the necessary twist, the length-one
quotient is the sign representation of $S_3$, while
$(\F_3^3/C)^*$ is its augmentation module. The dual of a
two-dimensional module is its determinant-inverse twist; here the
determinant is the sign of the $S_3$ factor. Multiplying the scalar
$C_2$ coordinate of $R$ by this sign is an automorphism of $R$
and supplies both asserted abstract identifications. Consequently
the same three displayed identities hold with $Q_*$ in place of
$Q$ and these new onto tails.
Both comparators have degree six, physical $h(9)=8$, and the
largest tail slope $(\log_2 3)/2$ is less than one. Apply
Theorem~\ref{thm:finite-comparator}.
\end{proof}

\begin{theorem}[All bounded symmetric-three counts]\label{thm:bounded-s3}
Every actual action with exact local natural $S_3$ and
$s\in\{2,3,4,6,8,12,18\}$ blocks has a complete exponentially
small forward row and a quadratically small scalar. Together with
Theorem~\ref{thm:prime-s3-capacity}, this covers every imprimitive
exact-local-$S_3$ action.
\end{theorem}
\begin{proof}
The $E=1$ case was disposed of above. When $D=1$, use
\eqref{eq:bounded-s3-layer} with comparator $T$ and the respective
capacities
\[
 \begin{array}{c|rrrrrr}s&3&4&6&8&12&18\\\hline
 H_3&1&1&2&1&3&7.\end{array}
\]
These follow from Lemma~\ref{lem:refined-capacity}; at six use
the semiregular three-cycle argument of Proposition~\ref{prop:bounded-c3},
lifting through a binary kernel by a Sylow three-subgroup. Every
margin $(h(3s)-s)/8-\gamma H_3$ is positive. At $s=2$, $D=1$
contradicts the exact local $S_3$ image; Proposition~\ref{prop:two-s3}
covers that count completely.

At $s=8$, irrespective of $D$, the prime-to-three bound is $H_3=1$;
retain $R$ in degree sixteen, with margin $1-\gamma>0$.
At $s=18$, a soluble transitive subgroup of $T$ gives $H_3\leq5$
and $H_2\leq2$, so the degree-eighteen comparator $T$ has margin
$7/2-5\gamma>0$. Otherwise use the actual semiregular $C_3$
witness for every nonsoluble degree-eighteen action recorded in the
finite input to Theorem~\ref{thm:three-twentieths}. Its lift through
$D$ has order three. The signed ternary module and its dual
untwist to six regular modules over $\F_3[C_3]$, so $H_3\leq6$.
The same comparator has margin $7/2-6\gamma>0$.

Now assume $D\ne1$ and use the actual isotype count $q$.
At $s=12$, if $q\leq2$, then $d=q\leq2$ and $H_3\leq3$.
For each literal normal quotient of $R$, a fixed map to its top
quotient has at most $2^{db/2}$ restrictions into its binary
kernel, times a fixed cocycle constant. Thus
$Z_J(R)\leq C_R2^{db/2}Z_J(T)$, including all normal graphs.
The margin is at least $2-3\gamma>0$.
If $q\geq3$, its isotype size is $4,3,2$ or one;
Lemma~\ref{lem:signed-isotypes} gives $H_3=1$ and the
degree-twenty-four comparator $R$ has margin $3/2-\gamma>0$.
At $s=6$, $q\geq2$ similarly gives isotype sizes $3,2,1$ and
margin $3/4-\gamma>0$ with degree-twelve comparator $R$.

At $s=4,6$, the remaining $q=1$ means that $D$ is central of
order two in $R$. Combine Lemma~\ref{lem:central-comparator}
with the capacities $H_3=1,2$. The comparator $T\times C_2$
has degrees six and eight; its margins are respectively
$3/4-\gamma>0$ and $5/4-2\gamma>0$.
The $s=4,q=2$ branch is Proposition~\ref{prop:four-s3-two-signs}.
For $s=4,q=4$ or $s=3,q=3$, every coordinate sign space is
distinct. Transitivity makes $E$ the whole ambient module,
irreducible and faithful for $R$. Since $R$ is soluble, the
common-derived estimate Lemma~\ref{lem:derived-affine} gives
$Z_J(U)\leq Z_J(R)+C_U2^{\gamma b}$; the comparator degrees
$2s<h(3s)$ and tail slopes meet Theorem~\ref{thm:finite-comparator}.
The only remaining case is $s=3,q=1$, covered by
Proposition~\ref{prop:three-s3-one-sign}.

The listed cases exhaust all divisors $q\mid s$ and the zero
kernels. Each multiplicative bound has strictly positive margin,
so Lemma~\ref{lem:finite-multiplicative} applies. There are finitely
many actual actions in these widths. Choose one block system and
the first applicable certificate for each action, then one
any-hot/cold partition on the complete accepted pointing union.
The cold coefficients are summed with their original
$\Delta(n,n-w)/a(U)$, and the hot sums with their original
$((n-w)!a(U)p_nG_{\lfloor n/2\rfloor})^{-1}$.
The finite sum preserves exponential and quadratic decay.
There is no additional action weight for isotypes, quotient
representations, or alternative block systems.
\end{proof}
